\documentclass[10pt,landscape]{article}
\usepackage[utf8]{inputenc}
\usepackage[margin=0.4in]{geometry}
\usepackage{amsmath,amssymb}
\usepackage{multicol}
\usepackage[breakable]{tcolorbox}
\usepackage{enumitem}
\usepackage{xcolor}

% Define custom colors for each topic
\definecolor{regression}{RGB}{230,240,255}
\definecolor{classification}{RGB}{255,240,230}
\definecolor{neural}{RGB}{240,255,240}
\definecolor{svm}{RGB}{255,240,255}
\definecolor{clustering}{RGB}{245,245,220}
\definecolor{mixture}{RGB}{230,255,255}
\definecolor{dimension}{RGB}{255,230,230}
\definecolor{graphical}{RGB}{240,240,255}
\definecolor{hmm}{RGB}{255,255,230}
\definecolor{mdp}{RGB}{255,235,245}
\definecolor{rl}{RGB}{235,255,235}

\pagestyle{empty}
\setlength{\columnsep}{0.2in}
\renewcommand{\arraystretch}{1.1}

\begin{document}
\footnotesize

\begin{multicols}{3}

%%%%%%%%%%%%%%%%%%%%%%%%%%%%%%%%%%%%%%%%%%
\begin{tcolorbox}[breakable, title=1. Regression Concepts, colback=regression]
\textbf{Linear Regression Overview:}
\begin{itemize}[leftmargin=*,noitemsep,topsep=0pt]
    \item Maps input $x \in \mathbb{R}^D$ to target $y \in \mathbb{R}$ using weights $w$
    \item Basic form: $y = w_0 + w_1x_1 + ... + w_Dx_D$
    \item Uses least squares loss to find optimal weights
    \item "Bias trick": Add $x_0=1$ to make notation cleaner: $y = w^T x$
\end{itemize}

\textbf{Key Considerations:}
\begin{itemize}[leftmargin=*,noitemsep,topsep=0pt]
    \item Basis functions ($\phi(x)$): Transform inputs to capture non-linear relationships
    \item Regularization: Prevent overfitting by penalizing large weights
    \item Bias-variance tradeoff: Balance between underfitting and overfitting
\end{itemize}

\textbf{Regularization Types:}
\begin{itemize}[leftmargin=*,noitemsep,topsep=0pt]
    \item Ridge (L2): $\frac{\lambda}{2}w^Tw$ - Shrinks all weights toward zero
    \item Lasso (L1): $\lambda\|w\|_1$ - Creates sparse solutions (some weights = 0)
    \item Elastic Net: Combination of L1 and L2
\end{itemize}

\textbf{Bayesian View:}
\begin{itemize}[leftmargin=*,noitemsep,topsep=0pt]
    \item Model noise: $y_n \sim \mathcal{N}(w^Tx_n, \beta^{-1})$
    \item Prior on weights: $w \sim \mathcal{N}(0, S_0^{-1})$
    \item MAP estimation equivalent to ridge regression with $\lambda = \frac{s}{\beta}$
\end{itemize}
\end{tcolorbox}

%%%%%%%%%%%%%%%%%%%%%%%%%%%%%%%%%%%%%%%%%%
\begin{tcolorbox}[breakable, title=1.1 Ridge Regression Derivation, colback=regression]
\textbf{When to use:} Use ridge regression when you suspect multicollinearity in your features or need to prevent overfitting in high-dimensional data.

\textbf{Objective:} Find $w_{ridge}$ that minimizes OLS with L2 regularization.

1. Start with ridge loss function:
\begin{align}
L(w) = \frac{1}{2}\sum_{n=1}^{N}(y_n - w^T\phi_n)^2 + \frac{\lambda}{2}w^Tw
\end{align}

2. In matrix form with design matrix $\Phi$ and target $y$:
\begin{align}
L(w) = \frac{1}{2}(y - \Phi w)^T(y - \Phi w) + \frac{\lambda}{2}w^Tw
\end{align}

3. Take gradient, set to zero:
\begin{align}
\nabla L(w) &= -\Phi^Ty + \Phi^T\Phi w + \lambda w = 0\\
\Rightarrow (\Phi^T\Phi + \lambda I)w &= \Phi^Ty
\end{align}

4. Optimal weights: $w_{ridge} = (\Phi^T\Phi + \lambda I)^{-1}\Phi^Ty$

\textbf{Comparison with OLS:} 
OLS solution is $w^* = (\Phi^T\Phi)^{-1}\Phi^Ty$. 
Ridge adds $\lambda I$ to $\Phi^T\Phi$, which:
\begin{itemize}[leftmargin=*,noitemsep,topsep=0pt]
    \item Makes matrix inversion more stable
    \item Shrinks weights toward zero
    \item Reduces overfitting by introducing bias
\end{itemize}
\end{tcolorbox}

%%%%%%%%%%%%%%%%%%%%%%%%%%%%%%%%%%%%%%%%%%
\begin{tcolorbox}[breakable, title=1.2 Maximum Likelihood Estimation, colback=regression]
\textbf{When to use:} Use MLE when you want to find parameters that make your observed data most probable under your chosen model.

\textbf{General MLE Procedure:}

1. Write likelihood function (probability of data given parameters):
$p(D|\theta) = \prod_{n=1}^{N} p(x_n|\theta)$

2. Take logarithm for easier optimization:
$\log p(D|\theta) = \sum_{n=1}^{N} \log p(x_n|\theta)$

3. Find parameters maximizing log-likelihood:
$\frac{\partial}{\partial \theta}\log p(D|\theta) = 0$

\textbf{Example: Linear Regression MLE}

1. Assume $y_n \sim \mathcal{N}(w^T x_n, \beta^{-1})$

2. Log-likelihood:
$\log p(y|X, w, \beta) = \frac{N}{2}\log\beta - \frac{N}{2}\log(2\pi) - \frac{\beta}{2}\sum_{n}(y_n - w^Tx_n)^2$

3. Take derivative, set to zero:
$\frac{\partial \log p}{\partial w} = \beta \sum_{n}(y_n - w^Tx_n)x_n = 0$

4. Solve with matrix notation:
$X^Ty = X^TXw \Rightarrow w^* = (X^TX)^{-1}X^Ty$

\textbf{Key insight:} MLE for linear regression with Gaussian noise is equivalent to minimizing the least squares loss!
\end{tcolorbox}

%%%%%%%%%%%%%%%%%%%%%%%%%%%%%%%%%%%%%%%%%%
\begin{tcolorbox}[breakable, title=1.3 Maximum A Posteriori Estimation, colback=regression]
\textbf{When to use:} Use MAP when you have prior knowledge about parameters and want to incorporate this information into your estimation.

\textbf{MAP General Procedure:}

1. Write posterior using Bayes' rule:
$p(\theta|D) \propto p(D|\theta)p(\theta)$

2. Take logarithm:
$\log p(\theta|D) = \log p(D|\theta) + \log p(\theta) + const.$

3. Maximize log posterior:
$\frac{\partial}{\partial \theta}[\log p(D|\theta) + \log p(\theta)] = 0$

\textbf{Example: Bayesian Linear Regression}

1. Gaussian prior: $w \sim \mathcal{N}(0, S_0^{-1})$

2. Log posterior:
$\log p(w|X,y,\beta) \propto -\frac{\beta}{2}\sum_n(y_n - w^Tx_n)^2 - \frac{1}{2}w^TS_0w$

3. Take derivative, set to zero:
$\beta\sum_n(y_n - w^Tx_n)x_n - S_0w = 0$

4. Matrix form solution:
$w_{MAP} = (\beta X^TX + S_0)^{-1}\beta X^Ty$

\textbf{Note:} With $S_0 = \lambda I$, we get ridge regression: 
$w_{MAP} = (X^TX + \frac{\lambda}{\beta}I)^{-1}X^Ty$
\end{tcolorbox}

%%%%%%%%%%%%%%%%%%%%%%%%%%%%%%%%%%%%%%%%%%
\begin{tcolorbox}[breakable, title=1.4 Posterior Parameters \& Predictive Distribution, colback=regression]
\textbf{When to use:} Use posterior predictive distribution when you want to make predictions that account for uncertainty in your parameters.

\textbf{Finding Posterior Parameters (Conjugate Distributions):}

1. Identify prior distribution and parameters
2. Identify likelihood function
3. Use conjugacy to determine posterior form
4. Derive posterior parameters by comparing coefficients

\textbf{Example: Gaussian-Gaussian conjugacy}
\begin{itemize}[leftmargin=*,noitemsep,topsep=0pt]
    \item Prior: $w \sim \mathcal{N}(\mu_0, \Sigma_0)$
    \item Likelihood: $y|x, w \sim \mathcal{N}(w^Tx, \sigma^2)$
    \item Posterior: $w|D \sim \mathcal{N}(\mu_N, \Sigma_N)$ where:
    \begin{align}
    \Sigma_N^{-1} &= \Sigma_0^{-1} + \frac{1}{\sigma^2}\sum_{n}x_nx_n^T\\
    \mu_N &= \Sigma_N(\Sigma_0^{-1}\mu_0 + \frac{1}{\sigma^2}\sum_{n}y_nx_n)
    \end{align}
\end{itemize}

\textbf{Posterior Predictive Distribution:}
\begin{align}
p(y^*|x^*, D) = \int p(y^*|x^*, \theta)p(\theta|D)d\theta
\end{align}

\textbf{For Gaussian posterior:}
$y^*|x^*, D \sim \mathcal{N}(\mu_N^Tx^*, \sigma^2 + x^{*T}\Sigma_Nx^*)$

\textbf{Key insight:} The posterior predictive variance includes both noise variance $\sigma^2$ and parameter uncertainty $x^{*T}\Sigma_Nx^*$.
\end{tcolorbox}

%%%%%%%%%%%%%%%%%%%%%%%%%%%%%%%%%%%%%%%%%%
\begin{tcolorbox}[breakable, title=2. Classification Concepts, colback=classification]
\textbf{Classification Overview:}
\begin{itemize}[leftmargin=*,noitemsep,topsep=0pt]
    \item Maps input $x$ to discrete class labels $y$
    \item Can be binary (two classes) or multi-class
    \item Decision boundaries separate different classes
\end{itemize}

\textbf{Common Classification Methods:}
\begin{itemize}[leftmargin=*,noitemsep,topsep=0pt]
    \item Discriminant functions: Direct mapping from inputs to class decisions
    \item Probabilistic discriminative models: Model $p(y|x)$ directly (logistic regression)
    \item Probabilistic generative models: Model $p(x,y)$ together (Naive Bayes)
\end{itemize}

\textbf{Loss Functions:}
\begin{itemize}[leftmargin=*,noitemsep,topsep=0pt]
    \item 0/1 Loss: Simple but non-differentiable, counts misclassifications
    \item Logistic Loss: Smooth approximation of 0/1 loss, used in logistic regression
    \item Hinge Loss: Used in SVMs, penalizes points close to or on wrong side of margin
\end{itemize}

\textbf{Training Techniques:}
\begin{itemize}[leftmargin=*,noitemsep,topsep=0pt]
    \item Gradient Descent: Iterative optimization to minimize loss
    \item Batch GD: Updates using all training examples
    \item SGD: Updates using one example at a time, faster but noisier
\end{itemize}

\textbf{Generative vs. Discriminative:}
\begin{itemize}[leftmargin=*,noitemsep,topsep=0pt]
    \item Discriminative: Models decision boundary directly
    \item Generative: Models how data is generated, can generate new samples
    \item Naive Bayes: Assumes feature independence given class
\end{itemize}
\end{tcolorbox}

%%%%%%%%%%%%%%%%%%%%%%%%%%%%%%%%%%%%%%%%%%
\begin{tcolorbox}[breakable, title=2.1 Logistic Regression Gradients, colback=classification]
\textbf{When to use:} Use logistic regression when you need probabilistic class membership for binary or multi-class classification with a linear decision boundary.

\textbf{Binary Logistic Regression:}

1. Model: $p(y=1|x) = \sigma(w^Tx) = \frac{1}{1+e^{-w^Tx}}$

2. Likelihood: $p(\{y_i\}_{i=1}^{N}|w) = \prod_{i=1}^{N}\hat{y}_i^{y_i}(1-\hat{y}_i)^{1-y_i}$
where $\hat{y}_i = \sigma(w^Tx_i)$

3. Negative log-likelihood (cross-entropy):
$E(w) = -\sum_{i=1}^{N}\{y_i\log\hat{y}_i + (1-y_i)\log(1-\hat{y}_i)\}$

4. Gradient: $\nabla E(w) = \sum_{i=1}^{N}(\hat{y}_i - y_i)x_i$

\textbf{Multi-class with Softmax:}
$\nabla_{w_j}E(W) = \sum_{i=1}^{N}(\hat{y}_{ij} - y_{ij})x_i$
where $\hat{y}_{ij} = \text{softmax}_j(Wx_i)$

\textbf{Softmax function:} $\text{softmax}_k(z) = \frac{\exp(z_k)}{\sum_{j=1}^K\exp(z_j)}$
\end{tcolorbox}

%%%%%%%%%%%%%%%%%%%%%%%%%%%%%%%%%%%%%%%%%%
\begin{tcolorbox}[breakable, title=2.2 Generative Models for Classification, colback=classification]
\textbf{When to use:} Use generative models when you want to model the data distribution, generate new samples, or when you have prior knowledge about how the data is generated.

\textbf{Framework:}
\begin{enumerate}[leftmargin=*,noitemsep,topsep=0pt]
    \item Define class prior $p(C_k)$ and class-conditional $p(x|C_k)$
    \item Joint distribution: $p(x, C_k) = p(C_k)p(x|C_k)$
    \item Full log-likelihood for binary classification:
    $\log p(\pi, \mu_1, \mu_2, \Sigma) = \sum_{i=1}^{N}y_i\log[\pi\mathcal{N}(x_i|\mu_1, \Sigma)] + (1-y_i)\log[(1-\pi)\mathcal{N}(x_i|\mu_2, \Sigma)]$
    \item Maximize to find optimal parameters
\end{enumerate}

\textbf{MLE Solutions (Gaussian with shared covariance):}
\begin{align}
\pi &= \frac{N_1}{N}\\
\mu_k &= \frac{1}{N_k}\sum_{i:y_i=k}x_i\\
\Sigma &= \frac{1}{N}\sum_{i=1}^{N}[y_i(x_i - \mu_1)(x_i - \mu_1)^T \\
&\quad + (1-y_i)(x_i - \mu_2)(x_i - \mu_2)^T]
\end{align}

\textbf{Using Lagrangian Method:}
1. Identify constraints (e.g., $\sum_{k}\pi_k = 1$)
2. Form Lagrangian: 
   $\mathcal{L}(\theta, \lambda) = \log p(D|\theta) + \lambda(c - \sum_{k}\pi_k)$
3. Take derivatives w.r.t. parameters and Lagrange multipliers
4. Solve the resulting system
\end{tcolorbox}

%%%%%%%%%%%%%%%%%%%%%%%%%%%%%%%%%%%%%%%%%%
\begin{tcolorbox}[breakable, title=2.3 Naive Bayes, colback=classification]
\textbf{When to use:} Use Naive Bayes when features can reasonably be assumed to be conditionally independent given the class, or when you have limited training data.

\textbf{Key Idea:} Assume features are conditionally independent given the class:
\begin{align}
p(x|y=C_k) = \prod_{d=1}^D p(x_d|y=C_k)
\end{align}

\textbf{Classification Rule:}
\begin{align}
\hat{y} = \arg\max_k p(y=C_k|x) = \arg\max_k p(y=C_k)\prod_{d=1}^D p(x_d|y=C_k)
\end{align}

\textbf{Parameter Estimation:}
\begin{itemize}[leftmargin=*,noitemsep,topsep=0pt]
    \item Class priors: $p(y=C_k) = \frac{N_k}{N}$
    \item For discrete features: $p(x_d=v|y=C_k) = \frac{\text{count}(x_d=v, y=C_k)}{\text{count}(y=C_k)}$
    \item For continuous features: Often assume Gaussian $p(x_d|y=C_k) = \mathcal{N}(\mu_{dk}, \sigma_{dk}^2)$
\end{itemize}

\textbf{Advantages:}
\begin{itemize}[leftmargin=*,noitemsep,topsep=0pt]
    \item Fast training and prediction
    \item Works well with small datasets
    \item Handles missing values naturally
\end{itemize}

\textbf{Limitations:}
\begin{itemize}[leftmargin=*,noitemsep,topsep=0pt]
    \item Independence assumption often violated
    \item Cannot learn feature interactions
\end{itemize}
\end{tcolorbox}

%%%%%%%%%%%%%%%%%%%%%%%%%%%%%%%%%%%%%%%%%%
\begin{tcolorbox}[breakable, title=2.4 Gradient Descent Implementation, colback=classification]
\textbf{When to use:} Use gradient descent when direct analytical solutions aren't possible or practical, especially with non-convex or high-dimensional loss functions.

\textbf{Basic Gradient Descent:}
\begin{enumerate}[leftmargin=*,noitemsep,topsep=0pt]
    \item Initialize parameters $\theta^{(0)}$
    \item Choose learning rate $\eta$
    \item For iterations $t=0,1,2...$:
    \begin{itemize}[leftmargin=*,noitemsep]
        \item Compute gradient: $\nabla L(\theta^{(t)})$
        \item Update: $\theta^{(t+1)} = \theta^{(t)} - \eta \nabla L(\theta^{(t)})$
        \item Check convergence
    \end{itemize}
\end{enumerate}

\textbf{For Logistic Regression:}
$w^{(t+1)} = w^{(t)} - \eta \sum_{i=1}^{N}(\sigma(w^{(t)T}x_i) - y_i)x_i$

\textbf{Stochastic Gradient Descent (SGD):}
Update using one or a batch of samples:
$w^{(t+1)} = w^{(t)} - \eta \sum_{i \in \text{batch}}(\sigma(w^{(t)T}x_i) - y_i)x_i$

\textbf{Learning Rate Selection:}
\begin{itemize}[leftmargin=*,noitemsep,topsep=0pt]
    \item Too large: May diverge or oscillate
    \item Too small: Slow convergence
    \item Adaptive rates: Decrease over time or use algorithms like Adam, RMSprop
\end{itemize}

\textbf{Convergence Criteria:}
\begin{itemize}[leftmargin=*,noitemsep,topsep=0pt]
    \item Small parameter change: $\|\theta^{(t+1)} - \theta^{(t)}\| < \epsilon$
    \item Small gradient: $\|\nabla L(\theta^{(t)})\| < \epsilon$
    \item Maximum iterations reached
\end{itemize}
\end{tcolorbox}

%%%%%%%%%%%%%%%%%%%%%%%%%%%%%%%%%%%%%%%%%%
\begin{tcolorbox}[breakable, title=2.5 Classifier Performance Analysis, colback=classification]
\textbf{Parametric vs. Non-parametric Classifiers:}

\textbf{Parametric} (logistic regression, Naive Bayes):
\begin{itemize}[leftmargin=*,noitemsep,topsep=0pt]
    \item Stronger assumptions about data
    \item Need fewer data points
    \item May have high bias if assumptions wrong
    \item Generally faster inference
\end{itemize}

\textbf{Non-parametric} (k-NN, kernel methods):
\begin{itemize}[leftmargin=*,noitemsep,topsep=0pt]
    \item Fewer assumptions about data
    \item Need more training data
    \item More flexible decision boundaries
    \item Higher variance/overfitting risk
    \item Often slower inference
\end{itemize}

\textbf{Expected Performance by Dataset:}
\begin{itemize}[leftmargin=*,noitemsep,topsep=0pt]
    \item Small datasets: Parametric models better
    \item Large datasets: Non-parametric capture complex patterns
    \item High-dimensional: Parametric more robust
    \item Linear structure: Linear models excel
    \item Complex patterns: Non-parametric/kernel methods better
\end{itemize}

\textbf{Model Selection Considerations:}
\begin{itemize}[leftmargin=*,noitemsep,topsep=0pt]
    \item Cross-validation to estimate generalization performance
    \item Use validation set for hyperparameter tuning
    \item Consider computational constraints for training and inference
    \item Evaluate against baseline models
\end{itemize}
\end{tcolorbox}

%%%%%%%%%%%%%%%%%%%%%%%%%%%%%%%%%%%%%%%%%%
\begin{tcolorbox}[breakable, title=3. Neural Networks Concepts, colback=neural]
\textbf{Neural Network Overview:}
\begin{itemize}[leftmargin=*,noitemsep,topsep=0pt]
    \item Universal function approximators
    \item Composed of layers of interconnected nodes (neurons)
    \item Each neuron applies weighted sum followed by activation function
    \item Learn adaptive basis functions automatically
\end{itemize}

\textbf{Key Components:}
\begin{itemize}[leftmargin=*,noitemsep,topsep=0pt]
    \item Input layer: Receives data features
    \item Hidden layers: Extract and transform features
    \item Output layer: Produces predictions
    \item Activation functions: Introduce non-linearity (ReLU, sigmoid, tanh)
    \item Weights and biases: Parameters learned during training
\end{itemize}

\textbf{Forward Propagation:}
\begin{align}
a^{(l)}_j &= \sum_{i=1}^{M^{(l-1)}} w^{(l)}_{ji} z^{(l-1)}_i + b^{(l)}_j\\
z^{(l)}_j &= h(a^{(l)}_j)
\end{align}
where $h$ is the activation function.

\textbf{Training Process:}
\begin{itemize}[leftmargin=*,noitemsep,topsep=0pt]
    \item Define loss function appropriate for task
    \item Backpropagation to calculate gradients
    \item Update weights using optimization algorithm (SGD, Adam, etc.)
\end{itemize}

\textbf{Preventing Overfitting:}
\begin{itemize}[leftmargin=*,noitemsep,topsep=0pt]
    \item Regularization (L1, L2, dropout)
    \item Early stopping
    \item Data augmentation
    \item Cross-validation
\end{itemize}
\end{tcolorbox}

%%%%%%%%%%%%%%%%%%%%%%%%%%%%%%%%%%%%%%%%%%
\begin{tcolorbox}[breakable, title=3.1 Neural Network Loss Functions, colback=neural]
\textbf{When to use:} Choose your loss function based on the specific task type (regression vs. classification) and data characteristics (noise, outliers).

\textbf{Regression Tasks:}
\begin{itemize}[leftmargin=*,noitemsep,topsep=0pt]
    \item MSE: $L(w) = \frac{1}{2}\sum_{n=1}^{N}(y_n - f(x_n, w))^2$
    \item MAE: $L(w) = \sum_{n=1}^{N}|y_n - f(x_n, w)|$
    \item Huber Loss: Combines MSE and MAE, robust to outliers
\end{itemize}

\textbf{Binary Classification:}
\begin{itemize}[leftmargin=*,noitemsep,topsep=0pt]
    \item Binary Cross-Entropy: 
    $L(w) = -\sum_{n=1}^{N}[y_n\log(\hat{y}_n) + (1-y_n)\log(1-\hat{y}_n)]$
    \item Hinge Loss: $L(w) = \sum_{n=1}^{N}\max(0, 1-y_n f(x_n, w))$
\end{itemize}

\textbf{Multi-class Classification:}
\begin{itemize}[leftmargin=*,noitemsep,topsep=0pt]
    \item Categorical Cross-Entropy: 
    $L(w) = -\sum_{n=1}^{N}\sum_{k=1}^{K}y_{nk}\log(\hat{y}_{nk})$
    \item KL-Divergence: Measures difference between predicted and true distributions
\end{itemize}

\textbf{Selection Criteria:}
\begin{itemize}[leftmargin=*,noitemsep,topsep=0pt]
    \item Task type (regression vs. classification)
    \item Sensitivity to outliers
    \item Distribution assumptions about data
    \item Model architecture and optimization constraints
\end{itemize}
\end{tcolorbox}

%%%%%%%%%%%%%%%%%%%%%%%%%%%%%%%%%%%%%%%%%%
\begin{tcolorbox}[breakable, title=3.2 Neural Network Architecture Analysis, colback=neural]
\textbf{When to use:} Consider these architectural choices when designing neural networks to balance capacity, trainability, and generalization performance.

\textbf{Layer Count Effects:}
\begin{itemize}[leftmargin=*,noitemsep,topsep=0pt]
    \item More layers: Greater capacity for complex functions, vanishing gradient risk
    \item Fewer layers: Reduced capacity, easier training, potential underfitting
\end{itemize}

\textbf{Neuron Count Effects:}
\begin{itemize}[leftmargin=*,noitemsep,topsep=0pt]
    \item More neurons: Increased capacity, overfitting risk
    \item Fewer neurons: Better generalization with limited data
\end{itemize}

\textbf{Activation Function Effects:}
\begin{itemize}[leftmargin=*,noitemsep,topsep=0pt]
    \item ReLU: Fast convergence, sparse activations, dying neuron risk
    \item Sigmoid/Tanh: Bounded outputs, vanishing gradient risk
    \item Leaky ReLU: Addresses dying neuron problem
\end{itemize}

\textbf{Network Architecture Types:}
\begin{itemize}[leftmargin=*,noitemsep,topsep=0pt]
    \item Feedforward: Standard, fully connected layers
    \item CNN: Specialized for spatial data (images)
    \item RNN: Specialized for sequential data (text, time series)
    \item Transformer: Self-attention based for sequence data
\end{itemize}
\end{tcolorbox}

%%%%%%%%%%%%%%%%%%%%%%%%%%%%%%%%%%%%%%%%%%
\begin{tcolorbox}[breakable, title=3.3 Backpropagation Computation, colback=neural]
\textbf{When to use:} Backpropagation is the standard method for training neural networks, used to efficiently compute gradients for all parameters.

\textbf{Forward Pass:} Compute activations at each layer:
\begin{align}
a^{(l)}_j &= \sum_{i=1}^{M^{(l-1)}} w^{(l)}_{ji} z^{(l-1)}_i\\
z^{(l)}_j &= h(a^{(l)}_j)
\end{align}
where $h$ is the activation function.

\textbf{Backward Pass:}
\begin{enumerate}[leftmargin=*,noitemsep,topsep=0pt]
    \item Compute output layer error:
    $\delta^{(L)}_j = \frac{\partial L}{\partial a^{(L)}_j}$
    
    For MSE: $\delta^{(L)}_j = (z^{(L)}_j - y_j) \cdot h'(a^{(L)}_j)$
    
    \item Backpropagate error:
    $\delta^{(l)}_j = h'(a^{(l)}_j) \sum_{k=1}^{M^{(l+1)}} w^{(l+1)}_{kj} \delta^{(l+1)}_k$
    
    \item Compute weight gradients:
    $\frac{\partial L}{\partial w^{(l)}_{ji}} = \delta^{(l)}_j z^{(l-1)}_i$
\end{enumerate}

\textbf{Practical Steps for Manual Calculation:}
\begin{enumerate}[leftmargin=*,noitemsep,topsep=0pt]
    \item Draw network diagram with all connections
    \item Label all activations and weights
    \item Compute forward pass values
    \item Derive output layer error (derivative of loss)
    \item Propagate errors backward layer by layer
    \item Calculate all weight gradients
\end{enumerate}
\end{tcolorbox}

%%%%%%%%%%%%%%%%%%%%%%%%%%%%%%%%%%%%%%%%%%
\begin{tcolorbox}[breakable, title=4. Support Vector Machines Concepts, colback=svm]
\textbf{SVM Overview:}
\begin{itemize}[leftmargin=*,noitemsep,topsep=0pt]
    \item Maximum margin classifier
    \item Finds hyperplane that maximally separates classes
    \item Uses support vectors (points closest to decision boundary)
    \item Can handle linearly and non-linearly separable data
\end{itemize}

\textbf{Key Components:}
\begin{itemize}[leftmargin=*,noitemsep,topsep=0pt]
    \item Decision function: $f(x) = w^Tx + w_0$
    \item Margin: Distance from decision boundary to nearest data points
    \item Support vectors: Data points on margin boundaries
    \item Regularization parameter C: Controls trade-off between margin size and misclassification
\end{itemize}

\textbf{Types of SVMs:}
\begin{itemize}[leftmargin=*,noitemsep,topsep=0pt]
    \item Hard-margin: No misclassifications allowed, requires linearly separable data
    \item Soft-margin: Allows some misclassifications, more robust to outliers
    \item Kernel SVMs: Transform data to higher dimensions for non-linear boundaries
\end{itemize}

\textbf{Common Kernels:}
\begin{itemize}[leftmargin=*,noitemsep,topsep=0pt]
    \item Linear: $K(x,z) = x^Tz$
    \item Polynomial: $K(x,z) = (1 + x^Tz)^q$
    \item RBF (Gaussian): $K(x,z) = \exp(-\gamma\|x-z\|^2)$
\end{itemize}

\textbf{Comparison to Logistic Regression:}
\begin{itemize}[leftmargin=*,noitemsep,topsep=0pt]
    \item SVM: Maximizes margin, focuses on boundary points
    \item Logistic regression: Probabilistic outputs, influenced by all points
\end{itemize}
\end{tcolorbox}

%%%%%%%%%%%%%%%%%%%%%%%%%%%%%%%%%%%%%%%%%%
\begin{tcolorbox}[breakable, title=4.1 Maximum Margin Classification, colback=svm]
\textbf{When to use:} Use hard-margin SVM when your data is linearly separable and you want a decision boundary with maximum distance to all classes.

\textbf{Hard Margin Formulation:}
\begin{align}
\min_{w,w_0} \frac{1}{2}\|w\|^2 \\
\text{subject to}~y_i(w^Tx_i + w_0) \geq 1,~\forall i
\end{align}

\textbf{Interpretation:}
\begin{itemize}[leftmargin=*,noitemsep,topsep=0pt]
    \item Decision boundary: $w^Tx + w_0 = 0$
    \item Margin width = $\frac{2}{\|w\|}$
    \item Support vectors exactly satisfy $y_i(w^Tx_i + w_0) = 1$
\end{itemize}

\textbf{Soft Margin Formulation:}
\begin{align}
\min_{w,w_0,\xi} \frac{1}{2}\|w\|^2 + C\sum_{i=1}^N \xi_i \\
\text{subject to}~y_i(w^Tx_i + w_0) \geq 1 - \xi_i,~\xi_i \geq 0,~\forall i
\end{align}

\textbf{Slack Variables $\xi_i$:}
\begin{itemize}[leftmargin=*,noitemsep,topsep=0pt]
    \item $\xi_i = 0$: Point correctly classified beyond margin
    \item $0 < \xi_i \leq 1$: Point correctly classified but inside margin
    \item $\xi_i > 1$: Point misclassified
\end{itemize}

\textbf{Parameter C:}
\begin{itemize}[leftmargin=*,noitemsep,topsep=0pt]
    \item Large C: Less regularization, small margin, fewer misclassifications
    \item Small C: More regularization, large margin, more misclassifications
\end{itemize}
\end{tcolorbox}

%%%%%%%%%%%%%%%%%%%%%%%%%%%%%%%%%%%%%%%%%%
\begin{tcolorbox}[breakable, title=4.2 The Kernel Trick, colback=svm]
\textbf{When to use:} Use kernel SVMs when your data is not linearly separable in the original feature space but may be separable in a higher-dimensional space.

\textbf{Key Insight:} The SVM dual formulation only requires inner products between data points, which can be replaced with kernel functions.

\textbf{Dual Formulation:}
\begin{align}
\max_{\alpha} \sum_{i=1}^N \alpha_i - \frac{1}{2}\sum_{i=1}^N\sum_{j=1}^N \alpha_i\alpha_j y_i y_j x_i^T x_j \\
\text{subject to}~\sum_{i=1}^N \alpha_i y_i = 0,~\alpha_i \geq 0,~\forall i
\end{align}

\textbf{Decision Function:}
\begin{align}
f(x) = \sum_{i=1}^N \alpha_i y_i K(x_i, x) + w_0
\end{align}

\textbf{Common Kernels:}
\begin{itemize}[leftmargin=*,noitemsep,topsep=0pt]
    \item Linear: $K(x,z) = x^Tz$
    \item Polynomial: $K(x,z) = (1 + x^Tz)^q$
    \begin{itemize}[leftmargin=*,noitemsep]
        \item For $q=2$, includes all constant, linear, and quadratic terms
    \end{itemize}
    \item RBF: $K(x,z) = \exp(-\gamma\|x-z\|^2)$
    \begin{itemize}[leftmargin=*,noitemsep]
        \item Corresponds to infinite-dimensional feature space
        \item Parameter $\gamma$ controls width of Gaussian
    \end{itemize}
\end{itemize}

\textbf{Mercer's Condition:} A valid kernel function must produce a positive semi-definite Gram matrix.

\textbf{Kernel Construction Rules:}
\begin{itemize}[leftmargin=*,noitemsep,topsep=0pt]
    \item Sum: $K_1 + K_2$
    \item Product: $K_1 \cdot K_2$
    \item Scalar: $c \cdot K$
    \item Exponentiation: $\exp(K)$
\end{itemize}
\end{tcolorbox}

%%%%%%%%%%%%%%%%%%%%%%%%%%%%%%%%%%%%%%%%%%
\begin{tcolorbox}[breakable, title=5. Clustering Concepts, colback=clustering]
\textbf{Clustering Overview:}
\begin{itemize}[leftmargin=*,noitemsep,topsep=0pt]
    \item Unsupervised learning technique to group similar data points
    \item No target labels provided
    \item Groups determined by similarity or distance metrics
    \item Used for data exploration, pattern discovery, and compression
\end{itemize}

\textbf{Key Types of Clustering:}
\begin{itemize}[leftmargin=*,noitemsep,topsep=0pt]
    \item Partitional: Divides data into non-overlapping groups (K-means)
    \item Hierarchical: Creates nested clusters (HAC)
    \item Density-based: Forms clusters based on dense regions (DBSCAN)
    \item Model-based: Assumes data comes from a mixture of distributions (GMM)
\end{itemize}

\textbf{Clustering Applications:}
\begin{itemize}[leftmargin=*,noitemsep,topsep=0pt]
    \item Customer segmentation
    \item Document categorization
    \item Image segmentation
    \item Anomaly detection
    \item Data preprocessing for supervised learning
\end{itemize}

\textbf{Evaluation Metrics:}
\begin{itemize}[leftmargin=*,noitemsep,topsep=0pt]
    \item Internal: Silhouette coefficient, Davies-Bouldin index
    \item External: Requires ground truth labels (Rand index, F-measure)
\end{itemize}
\end{tcolorbox}

%%%%%%%%%%%%%%%%%%%%%%%%%%%%%%%%%%%%%%%%%%
\begin{tcolorbox}[breakable, title=5.1 K-means Clustering, colback=clustering]
\textbf{When to use:} Use K-means when you need spherical clusters of similar sizes and the number of clusters is known in advance.

\textbf{K-means Objective:}
\begin{align}
L(\{r_n\}, \{\mu_c\}) = \sum_{n=1}^N\sum_{c=1}^C r_{nc}\|x_n - \mu_c\|^2
\end{align}
where $r_{nc} = 1$ if point $x_n$ belongs to cluster $c$, otherwise 0.

\textbf{Lloyd's Algorithm:}
\begin{enumerate}[leftmargin=*,noitemsep,topsep=0pt]
    \item Initialize cluster centers $\mu_1, \mu_2, ..., \mu_C$
    \item Repeat until convergence:
    \begin{itemize}[leftmargin=*,noitemsep]
        \item Assignment: Assign each point to nearest cluster
        $r_{nc} = 1$ if $c = \arg\min_{c'} \|x_n - \mu_{c'}\|^2$, else 0
        \item Update: Recalculate cluster centers
        $\mu_c = \frac{\sum_{n=1}^N r_{nc}x_n}{\sum_{n=1}^N r_{nc}}$
    \end{itemize}
\end{enumerate}

\textbf{Choosing Number of Clusters:}
\begin{itemize}[leftmargin=*,noitemsep,topsep=0pt]
    \item Elbow method: Plot loss vs. number of clusters
    \item Silhouette analysis: Measure cluster cohesion and separation
    \item Gap statistic: Compare with random reference distribution
\end{itemize}

\textbf{K-means++:} Improved initialization
\begin{enumerate}[leftmargin=*,noitemsep,topsep=0pt]
    \item Choose first center uniformly at random
    \item For each subsequent center, select points with probability proportional to squared distance from nearest existing center
\end{enumerate}

\textbf{K-medoids:} Uses actual data points as centers
\begin{itemize}[leftmargin=*,noitemsep,topsep=0pt]
    \item More robust to outliers
    \item Works with arbitrary distance metrics
    \item Higher computational cost
\end{itemize}
\end{tcolorbox}

%%%%%%%%%%%%%%%%%%%%%%%%%%%%%%%%%%%%%%%%%%
\begin{tcolorbox}[breakable, title=5.2 Hierarchical Agglomerative Clustering, colback=clustering]
\textbf{When to use:} Use HAC when you want to discover hierarchical structure in your data or when the number of clusters is not known in advance.

\textbf{HAC Algorithm:}
\begin{enumerate}[leftmargin=*,noitemsep,topsep=0pt]
    \item Start with each point in its own cluster
    \item Repeat until only one cluster remains:
    \begin{itemize}[leftmargin=*,noitemsep]
        \item Find closest pair of clusters using linkage criterion
        \item Merge these clusters
        \item Update distances to new cluster
    \end{itemize}
    \item Result is a dendrogram showing merge history
\end{enumerate}

\textbf{Linkage Criteria:}
\begin{itemize}[leftmargin=*,noitemsep,topsep=0pt]
    \item Single (min): $d(C_i, C_j) = \min_{x \in C_i, y \in C_j} \|x - y\|$
    \begin{itemize}[leftmargin=*,noitemsep]
        \item Tends to create elongated, chained clusters
    \end{itemize}
    \item Complete (max): $d(C_i, C_j) = \max_{x \in C_i, y \in C_j} \|x - y\|$
    \begin{itemize}[leftmargin=*,noitemsep]
        \item Creates compact, similar-sized clusters
    \end{itemize}
    \item Average: $d(C_i, C_j) = \frac{1}{|C_i||C_j|}\sum_{x \in C_i}\sum_{y \in C_j}\|x - y\|$
    \begin{itemize}[leftmargin=*,noitemsep]
        \item Compromise between single and complete
    \end{itemize}
    \item Ward: Minimizes increase in variance after merging
\end{itemize}

\textbf{Advantages over K-means:}
\begin{itemize}[leftmargin=*,noitemsep,topsep=0pt]
    \item No need to specify number of clusters in advance
    \item Can create clusters of any shape
    \item Provides hierarchical relationship view
    \item Deterministic (no random initialization)
\end{itemize}

\textbf{Limitations:}
\begin{itemize}[leftmargin=*,noitemsep,topsep=0pt]
    \item O(N²) space complexity for distance matrix
    \item O(N³) time complexity in worst case
    \item Cannot handle very large datasets
\end{itemize}
\end{tcolorbox}

%%%%%%%%%%%%%%%%%%%%%%%%%%%%%%%%%%%%%%%%%%
\begin{tcolorbox}[breakable, title=6. Mixture Models Concepts, colback=mixture]
\textbf{Mixture Models Overview:}
\begin{itemize}[leftmargin=*,noitemsep,topsep=0pt]
    \item Probabilistic models for data from multiple sources/distributions
    \item Each data point belongs to a latent (unobserved) class
    \item Models joint distribution p(x,z) of data x and latent class z
    \item Can be used for clustering, density estimation, and generative modeling
\end{itemize}

\textbf{General Form:}
\begin{align}
p(x) = \sum_{k=1}^K \pi_k p(x|z=k)
\end{align}
where $\pi_k$ are mixing coefficients ($\sum_k \pi_k = 1$) and $p(x|z=k)$ are component distributions.

\textbf{Estimation Challenge:}
\begin{itemize}[leftmargin=*,noitemsep,topsep=0pt]
    \item Maximum likelihood: $\log p(X) = \sum_{i=1}^N \log \sum_{k=1}^K \pi_k p(x_i|z=k)$
    \item Sum inside logarithm prevents direct optimization
    \item Solution: Use Expectation-Maximization (EM) algorithm
\end{itemize}

\textbf{Types of Mixture Models:}
\begin{itemize}[leftmargin=*,noitemsep,topsep=0pt]
    \item Gaussian Mixture Models (GMMs): Component distributions are Gaussian
    \item Mixture of Multinomials: For discrete data (text, genomics)
    \item Mixture of Experts: Different predictive models for different regions
    \item Latent Dirichlet Allocation: Documents as mixtures of topics
\end{itemize}
\end{tcolorbox}

%%%%%%%%%%%%%%%%%%%%%%%%%%%%%%%%%%%%%%%%%%
\begin{tcolorbox}[breakable, title=6.1 Expectation-Maximization (EM) Algorithm, colback=mixture]
\textbf{When to use:} Use EM when you have latent variable models where direct maximum likelihood is intractable, especially for mixture models.

\textbf{General Procedure:}
\begin{enumerate}[leftmargin=*,noitemsep,topsep=0pt]
    \item Initialize parameters $\theta^{(0)}$
    \item Repeat until convergence:
    \begin{itemize}[leftmargin=*,noitemsep]
        \item E-step: Compute 
        $Q(\theta|\theta^{(t)}) = \mathbb{E}_{Z|X,\theta^{(t)}}[\log p(X, Z|\theta)]$
        \item M-step: $\theta^{(t+1)} = \arg\max_\theta Q(\theta|\theta^{(t)})$
    \end{itemize}
\end{enumerate}

\textbf{Theoretical Foundation:}
\begin{itemize}[leftmargin=*,noitemsep,topsep=0pt]
    \item EM maximizes a lower bound (ELBO) on the log likelihood:
    \begin{align}
    \log p(X|\theta) \geq \mathbb{E}_{Z|X,\theta^{(t)}}[\log p(X,Z|\theta)] - \mathbb{E}_{Z|X,\theta^{(t)}}[\log p(Z|X,\theta^{(t)})]
    \end{align}
    \item E-step: Compute posterior distribution of latent variables
    \item M-step: Maximize expected complete-data log likelihood
    \item Guaranteed to increase log likelihood at each iteration
\end{itemize}

\textbf{Mathematical Properties:}
\begin{itemize}[leftmargin=*,noitemsep,topsep=0pt]
    \item Converges to local maximum of likelihood
    \item ELBO becomes tight when using true posterior
    \item Different initializations may lead to different local maxima
\end{itemize}

\textbf{Applications Beyond Mixture Models:}
\begin{itemize}[leftmargin=*,noitemsep,topsep=0pt]
    \item Hidden Markov Models
    \item Latent factor models (PCA, factor analysis)
    \item Missing data imputation
\end{itemize}
\end{tcolorbox}

%%%%%%%%%%%%%%%%%%%%%%%%%%%%%%%%%%%%%%%%%%
\begin{tcolorbox}[breakable, title=6.2 Gaussian Mixture Models (GMM), colback=mixture]
\textbf{When to use:} Use GMMs when your data appears to come from multiple Gaussian distributions or when you need soft clustering with uncertainty estimates.

\textbf{Model:} $p(x) = \sum_{k=1}^K \pi_k \mathcal{N}(x|\mu_k, \Sigma_k)$

\textbf{EM Algorithm for GMM:}
\begin{enumerate}[leftmargin=*,noitemsep,topsep=0pt]
    \item E-step: Calculate responsibilities (posterior probabilities)
    \begin{align}
    \gamma_{ik} = \frac{\pi_k \mathcal{N}(x_i|\mu_k, \Sigma_k)}{\sum_{j=1}^K \pi_j \mathcal{N}(x_i|\mu_j, \Sigma_j)}
    \end{align}
    
    \item M-step: Update parameters
    \begin{align}
    N_k &= \sum_{i=1}^N \gamma_{ik}\\
    \pi_k &= \frac{N_k}{N}\\
    \mu_k &= \frac{1}{N_k}\sum_{i=1}^N \gamma_{ik}x_i\\
    \Sigma_k &= \frac{1}{N_k}\sum_{i=1}^N \gamma_{ik}(x_i - \mu_k)(x_i - \mu_k)^T
    \end{align}
\end{enumerate}

\textbf{GMM vs. K-means:}
\begin{itemize}[leftmargin=*,noitemsep,topsep=0pt]
    \item GMM provides soft assignments (probabilities)
    \item GMM models covariance structure
    \item GMM is more flexible but requires more parameters
    \item K-means is a special case of GMM with:
    \begin{itemize}[leftmargin=*,noitemsep]
        \item Equal mixing coefficients $\pi_k = 1/K$
        \item Spherical covariances $\Sigma_k = \epsilon I$
        \item $\epsilon \to 0$ (hard assignments)
    \end{itemize}
\end{itemize}

\textbf{Choosing K:}
\begin{itemize}[leftmargin=*,noitemsep,topsep=0pt]
    \item BIC: $\text{BIC} = -2\log p(X|\hat{\theta}) + p\log N$
    \item AIC: $\text{AIC} = -2\log p(X|\hat{\theta}) + 2p$
    \item Cross-validation
\end{itemize}
\end{tcolorbox}

%%%%%%%%%%%%%%%%%%%%%%%%%%%%%%%%%%%%%%%%%%
\begin{tcolorbox}[breakable, title=6.3 Latent Dirichlet Allocation (LDA), colback=mixture]
\textbf{When to use:} Use LDA for topic modeling in text data when you want to discover latent topics and their distributions across documents.

\textbf{Key Idea:} Documents are mixtures of topics, and topics are distributions over words.

\textbf{Generative Process:}
\begin{enumerate}[leftmargin=*,noitemsep,topsep=0pt]
    \item For each document:
    \begin{itemize}[leftmargin=*,noitemsep]
        \item Draw topic mixture $\theta_d \sim \text{Dir}(\alpha)$
    \end{itemize}
    \item For each topic:
    \begin{itemize}[leftmargin=*,noitemsep]
        \item Draw word distribution $\phi_k \sim \text{Dir}(\beta)$
    \end{itemize}
    \item For each word position in each document:
    \begin{itemize}[leftmargin=*,noitemsep]
        \item Draw topic $z_{d,n} \sim \text{Cat}(\theta_d)$
        \item Draw word $w_{d,n} \sim \text{Cat}(\phi_{z_{d,n}})$
    \end{itemize}
\end{enumerate}

\textbf{LDA as an Admixture Model:}
\begin{itemize}[leftmargin=*,noitemsep,topsep=0pt]
    \item Documents can belong to multiple topics
    \item Words within a document can come from different topics
    \item Extends mixture models with hierarchical structure
\end{itemize}

\textbf{Inference Methods:}
\begin{itemize}[leftmargin=*,noitemsep,topsep=0pt]
    \item Variational inference
    \item Gibbs sampling
    \item Both approximate posterior distributions for latent variables
\end{itemize}

\textbf{LDA Applications:}
\begin{itemize}[leftmargin=*,noitemsep,topsep=0pt]
    \item Document clustering
    \item Information retrieval
    \item Recommendation systems
    \item Feature extraction for text classification
\end{itemize}
\end{tcolorbox}

%%%%%%%%%%%%%%%%%%%%%%%%%%%%%%%%%%%%%%%%%%
\begin{tcolorbox}[breakable, title=7. Dimensionality Reduction Concepts, colback=dimension]
\textbf{Dimensionality Reduction Overview:}
\begin{itemize}[leftmargin=*,noitemsep,topsep=0pt]
    \item Reduces high-dimensional data to lower dimensions
    \item Preserves important structure/information
    \item Used for visualization, compression, and preprocessing
    \item Helps combat curse of dimensionality
\end{itemize}

\textbf{Key Approaches:}
\begin{itemize}[leftmargin=*,noitemsep,topsep=0pt]
    \item Linear: PCA, LDA, Factor Analysis
    \item Non-linear: t-SNE, UMAP, Autoencoders
    \item Feature selection: Filter, wrapper, and embedded methods
\end{itemize}

\textbf{Common Applications:}
\begin{itemize}[leftmargin=*,noitemsep,topsep=0pt]
    \item Data visualization (2D/3D plots)
    \item Noise reduction
    \item Feature extraction
    \item Improving model performance
    \item Avoiding overfitting in high-dimensional spaces
\end{itemize}

\textbf{Evaluation Metrics:}
\begin{itemize}[leftmargin=*,noitemsep,topsep=0pt]
    \item Reconstruction error
    \item Variance explained
    \item Distance/neighborhood preservation
    \item Downstream task performance
\end{itemize}
\end{tcolorbox}

%%%%%%%%%%%%%%%%%%%%%%%%%%%%%%%%%%%%%%%%%%
\begin{tcolorbox}[breakable, title=7.1 Principal Component Analysis (PCA), colback=dimension]
\textbf{When to use:} Use PCA when you need linear dimensionality reduction that preserves maximum variance in your data, especially for visualization or preprocessing.

\textbf{PCA Steps:}
\begin{enumerate}[leftmargin=*,noitemsep,topsep=0pt]
    \item Center data: $\tilde{X} = X - \mu_X$
    \item Compute covariance: $S = \frac{1}{N}\tilde{X}^T\tilde{X}$
    \item Find eigenvalues/vectors: $S v_i = \lambda_i v_i$
    \item Sort eigenvectors by decreasing eigenvalue
    \item Select top $k$ eigenvectors as matrix $W_k$
    \item Project data: $Z = \tilde{X}W_k$
\end{enumerate}

\textbf{Two Key Perspectives:}
\begin{itemize}[leftmargin=*,noitemsep,topsep=0pt]
    \item Variance Maximization: Principal components are directions of maximum variance
    \item Reconstruction Error Minimization: Minimizes squared error when projecting back to original space
\end{itemize}

\textbf{Using SVD:}
\begin{itemize}[leftmargin=*,noitemsep,topsep=0pt]
    \item Decompose data matrix: $X = USV^T$
    \item Principal components = columns of $V$
    \item Variances along components = $\frac{s_i^2}{N}$
    \item Projected data: $Z = US$
\end{itemize}

\textbf{Choosing Components:}
\begin{itemize}[leftmargin=*,noitemsep,topsep=0pt]
    \item Scree plot: Plot eigenvalues and look for elbow
    \item Variance explained: $\frac{\sum_{i=1}^k \lambda_i}{\sum_{i=1}^D \lambda_i}$
    \item Kaiser criterion: Keep components with eigenvalues > 1
\end{itemize}

\textbf{Limitations:}
\begin{itemize}[leftmargin=*,noitemsep,topsep=0pt]
    \item Only captures linear relationships
    \item Sensitive to outliers
    \item Components may be hard to interpret
\end{itemize}
\end{tcolorbox}

%%%%%%%%%%%%%%%%%%%%%%%%%%%%%%%%%%%%%%%%%%
\begin{tcolorbox}[breakable, title=8. Graphical Models Concepts, colback=graphical]
\textbf{Graphical Models Overview:}
\begin{itemize}[leftmargin=*,noitemsep,topsep=0pt]
    \item Visual representation of probabilistic relationships
    \item Nodes represent random variables
    \item Edges represent dependence relationships
    \item Enables efficient inference and parameter learning
\end{itemize}

\textbf{Types of Graphical Models:}
\begin{itemize}[leftmargin=*,noitemsep,topsep=0pt]
    \item Directed (Bayesian Networks): Represent causal relationships
    \item Undirected (Markov Networks): Represent symmetric relationships
    \item Factor Graphs: General representation of factorized distributions
\end{itemize}

\textbf{Key Components:}
\begin{itemize}[leftmargin=*,noitemsep,topsep=0pt]
    \item Random variables: Observed (shaded) and latent (unshaded)
    \item Parameters: Model parameters shown as small dots
    \item Conditional probability tables (CPTs): Define p(node|parents)
    \item Plate notation: Represents repeated variables
\end{itemize}

\textbf{Joint Distribution Factorization:}
\begin{itemize}[leftmargin=*,noitemsep,topsep=0pt]
    \item Bayesian Network: $p(X_1, ..., X_n) = \prod_{i=1}^n p(X_i|\text{pa}(X_i))$
    \item Benefits: More compact representation, explicit independence assumptions
\end{itemize}
\end{tcolorbox}

%%%%%%%%%%%%%%%%%%%%%%%%%%%%%%%%%%%%%%%%%%
\begin{tcolorbox}[breakable, title=8.1 Bayesian Networks and Inference, colback=graphical]
\textbf{When to use:} Use Bayesian networks when you need to model probabilistic dependencies between variables, especially causal relationships.

\textbf{Constructing Bayesian Networks:}
\begin{enumerate}[leftmargin=*,noitemsep,topsep=0pt]
    \item For variable ordering $X_1, X_2, ..., X_n$:
    \begin{itemize}[leftmargin=*,noitemsep]
        \item Add $X_1$ (no edges)
        \item For each $X_i$, $i > 1$:
        \begin{itemize}[leftmargin=*,noitemsep]
            \item Add node $X_i$
            \item Add minimal edges from $\{X_1,...,X_{i-1}\}$ to $X_i$ to ensure conditional independence
        \end{itemize}
    \end{itemize}
\end{enumerate}

\textbf{Variable Elimination:}
\begin{enumerate}[leftmargin=*,noitemsep,topsep=0pt]
    \item Express joint as factor product: $p(X) = \prod_{i=1}^n p(X_i|\text{pa}(X_i))$
    \item Identify query $Q$ and evidence $E=e$
    \item Remove irrelevant variables and instantiate evidence
    \item Choose elimination order (leaves-first often efficient)
    \item For each variable $X$ in elimination order:
    \begin{itemize}[leftmargin=*,noitemsep]
        \item Collect all factors containing $X$
        \item Multiply these factors
        \item Sum out $X$ from resulting factor
        \item Replace original factors with new factor
    \end{itemize}
    \item Multiply remaining factors and normalize
\end{enumerate}

\textbf{D-Separation Test:}
\begin{enumerate}[leftmargin=*,noitemsep,topsep=0pt]
    \item Find all undirected paths between $X$ and $Y$
    \item Check if each path is blocked:
    \begin{itemize}[leftmargin=*,noitemsep]
        \item Chain ($A \rightarrow B \rightarrow C$): Blocked if $B \in Z$
        \item Fork ($A \leftarrow B \rightarrow C$): Blocked if $B \in Z$
        \item Collider ($A \rightarrow B \leftarrow C$): Blocked unless $B$ or any descendant of $B$ is in $Z$
    \end{itemize}
    \item $X$ and $Y$ are d-separated given $Z$ if all paths blocked
    \item Implication: If d-separated, then $X \perp Y | Z$
\end{enumerate}
\end{tcolorbox}

%%%%%%%%%%%%%%%%%%%%%%%%%%%%%%%%%%%%%%%%%%
\begin{tcolorbox}[breakable, title=9. Hidden Markov Models Concepts, colback=hmm]
\textbf{HMM Overview:}
\begin{itemize}[leftmargin=*,noitemsep,topsep=0pt]
    \item Models sequential data with hidden states
    \item States follow Markov property (depends only on previous state)
    \item Each state generates an observation (emission)
    \item Useful for temporal or sequential data
\end{itemize}

\textbf{Key Components:}
\begin{itemize}[leftmargin=*,noitemsep,topsep=0pt]
    \item Hidden states $s_t$: Unobserved state at time $t$
    \item Emissions $x_t$: Observed values at time $t$
    \item Initial state distribution: $p(s_1)$
    \item Transition probabilities: $p(s_t|s_{t-1})$
    \item Emission probabilities: $p(x_t|s_t)$
\end{itemize}

\textbf{Common Problems:}
\begin{itemize}[leftmargin=*,noitemsep,topsep=0pt]
    \item Inference: Determine p(states|observations)
    \item Learning: Estimate parameters from observations
    \item Decoding: Find most likely state sequence
\end{itemize}

\textbf{Common Applications:}
\begin{itemize}[leftmargin=*,noitemsep,topsep=0pt]
    \item Speech recognition
    \item Natural language processing
    \item Biological sequence analysis
    \item Financial time series
\end{itemize}
\end{tcolorbox}

%%%%%%%%%%%%%%%%%%%%%%%%%%%%%%%%%%%%%%%%%%
\begin{tcolorbox}[breakable, title=9.1 Forward-Backward Algorithm for HMMs, colback=hmm]
\textbf{When to use:} Use the Forward-Backward algorithm for efficient inference in HMMs, particularly for computing marginal probabilities of hidden states.

\textbf{Forward Pass ($\alpha$ values):}
\begin{align}
\alpha_t(s_t) &= p(x_1, ..., x_t, s_t)\\
\alpha_1(s_1) &= p(x_1|s_1)p(s_1)\\
\alpha_t(s_t) &= p(x_t|s_t) \sum_{s_{t-1}} p(s_t|s_{t-1})\alpha_{t-1}(s_{t-1})
\end{align}

\textbf{Backward Pass ($\beta$ values):}
\begin{align}
\beta_t(s_t) &= p(x_{t+1}, ..., x_n|s_t)\\
\beta_n(s_n) &= 1\\
\beta_t(s_t) &= \sum_{s_{t+1}} p(s_{t+1}|s_t)p(x_{t+1}|s_{t+1})\beta_{t+1}(s_{t+1})
\end{align}

\textbf{Inference with $\alpha$ and $\beta$:}
\begin{align}
p(x_1, ..., x_n) &= \sum_{s_t} \alpha_t(s_t)\beta_t(s_t)\\
p(s_t|x_1, ..., x_n) &\propto \alpha_t(s_t)\beta_t(s_t)\\
p(s_t, s_{t+1}|x_1, ..., x_n) &\propto \alpha_t(s_t)p(s_{t+1}|s_t)\\
&\quad\cdot p(x_{t+1}|s_{t+1})\beta_{t+1}(s_{t+1})
\end{align}

\textbf{Viterbi Algorithm (Best Path):}
\begin{enumerate}[leftmargin=*,noitemsep,topsep=0pt]
    \item Initialize:
    \begin{align}
    \delta_1(i) &= p(s_1 = i)p(x_1|s_1 = i)\\
    \psi_1(i) &= 0
    \end{align}
    
    \item Recursion:
    \begin{align}
    \delta_t(j) &= \max_i [\delta_{t-1}(i)p(s_t = j|s_{t-1} = i)]\\
    &\quad\cdot p(x_t|s_t = j)\\
    \psi_t(j) &= \arg\max_i [\delta_{t-1}(i)p(s_t = j|s_{t-1} = i)]
    \end{align}
    
    \item Termination:
    \begin{align}
    p^* &= \max_i [\delta_T(i)]\\
    s^*_T &= \arg\max_i [\delta_T(i)]
    \end{align}
    
    \item Backtracking:
    $s^*_t = \psi_{t+1}(s^*_{t+1})$ for $t = T-1,...,1$
\end{enumerate}

\textbf{HMM Training with EM:}
\begin{itemize}[leftmargin=*,noitemsep,topsep=0pt]
    \item E-step: Compute expected state occupancy and transitions using forward-backward
    \item M-step: Update parameters using these expectations
\end{itemize}
\end{tcolorbox}

%%%%%%%%%%%%%%%%%%%%%%%%%%%%%%%%%%%%%%%%%%
\begin{tcolorbox}[breakable, title=10. Markov Decision Processes Concepts, colback=mdp]
\textbf{MDP Overview:}
\begin{itemize}[leftmargin=*,noitemsep,topsep=0pt]
    \item Framework for decision-making under uncertainty
    \item Agent takes actions in states to maximize rewards
    \item Environment transitions between states based on actions
    \item Goal: Find optimal policy (mapping from states to actions)
\end{itemize}

\textbf{Key Components:}
\begin{itemize}[leftmargin=*,noitemsep,topsep=0pt]
    \item States $S$: Possible situations for the agent
    \item Actions $A$: Possible choices in each state
    \item Transition function $p(s'|s,a)$: Probability of next state
    \item Reward function $r(s,a)$: Immediate reward for action
    \item Discount factor $\gamma$: Value of future rewards
    \item Policy $\pi(a|s)$: Strategy for choosing actions
\end{itemize}

\textbf{Value Functions:}
\begin{itemize}[leftmargin=*,noitemsep,topsep=0pt]
    \item State value $V^\pi(s)$: Expected return starting from state $s$, following policy $\pi$
    \item Action value $Q^\pi(s,a)$: Expected return starting from state $s$, taking action $a$, then following policy $\pi$
\end{itemize}

\textbf{Planning Horizons:}
\begin{itemize}[leftmargin=*,noitemsep,topsep=0pt]
    \item Finite horizon: Fixed number of time steps
    \item Infinite horizon: Unbounded time with discounting
\end{itemize}
\end{tcolorbox}

%%%%%%%%%%%%%%%%%%%%%%%%%%%%%%%%%%%%%%%%%%
\begin{tcolorbox}[breakable, title=10.1 MDP Planning Algorithms, colback=mdp]
\textbf{When to use:} Use these planning algorithms when you have a fully defined MDP (states, actions, transitions, rewards) and need to find the optimal policy.

\textbf{Value Iteration:}
\begin{enumerate}[leftmargin=*,noitemsep,topsep=0pt]
    \item Initialize $V(s) = 0$ for all states
    \item Repeat until convergence:
    \begin{align}
    V'(s) &\leftarrow \max_a \left[r(s, a) + \gamma \sum_{s'} p(s'|s, a)V(s')\right]\\
    V(s) &\leftarrow V'(s)
    \end{align}
    \item Extract policy: 
    $\pi^*(s) = \arg\max_a \left[r(s, a) + \gamma \sum_{s'} p(s'|s, a)V^*(s')\right]$
\end{enumerate}

\textbf{Policy Iteration:}
\begin{enumerate}[leftmargin=*,noitemsep,topsep=0pt]
    \item Initialize policy $\pi$
    \item Repeat until convergence:
    \begin{itemize}[leftmargin=*,noitemsep]
        \item Policy Evaluation: Calculate $V^\pi(s)$
        \begin{align}
        V^\pi(s) = r(s, \pi(s)) + \gamma \sum_{s'} p(s'|s, \pi(s))V^\pi(s')
        \end{align}
        \item Policy Improvement:
        \begin{align}
        \pi'(s) &\leftarrow \arg\max_a \left[r(s, a) + \gamma \sum_{s'} p(s'|s, a)V^\pi(s')\right]\\
        \pi(s) &\leftarrow \pi'(s)
        \end{align}
    \end{itemize}
\end{enumerate}

\textbf{Value Iteration vs. Policy Iteration:}
\begin{itemize}[leftmargin=*,noitemsep,topsep=0pt]
    \item Value Iteration: More iterations, less computation per iteration
    \item Policy Iteration: Fewer iterations, more computation per iteration
    \item Value Iteration: O($|S|^2|A|$) per iteration
    \item Policy Evaluation: O($|S|^3$) for exact solution
\end{itemize}

\textbf{Bellman Equations:}
\begin{itemize}[leftmargin=*,noitemsep,topsep=0pt]
    \item Bellman Expectation Equation:
    $V^\pi(s) = \sum_a \pi(a|s)[r(s,a) + \gamma\sum_{s'} p(s'|s,a)V^\pi(s')]$
    \item Bellman Optimality Equation:
    $V^*(s) = \max_a [r(s,a) + \gamma\sum_{s'} p(s'|s,a)V^*(s')]$
\end{itemize}
\end{tcolorbox}

%%%%%%%%%%%%%%%%%%%%%%%%%%%%%%%%%%%%%%%%%%
\begin{tcolorbox}[breakable, title=11. Reinforcement Learning Concepts, colback=rl]
\textbf{RL Overview:}
\begin{itemize}[leftmargin=*,noitemsep,topsep=0pt]
    \item Learning from interaction with environment
    \item Agent takes actions, receives rewards, learns from experience
    \item No supervision, only feedback through rewards
    \item Goal: Learn policy that maximizes expected cumulative reward
\end{itemize}

\textbf{Key Components:}
\begin{itemize}[leftmargin=*,noitemsep,topsep=0pt]
    \item Agent: Decision-maker and learner
    \item Environment: Everything outside the agent
    \item State: Current situation
    \item Action: Agent's choice
    \item Reward: Feedback signal
    \item Policy: Strategy for choosing actions
\end{itemize}

\textbf{Exploration vs. Exploitation:}
\begin{itemize}[leftmargin=*,noitemsep,topsep=0pt]
    \item Exploration: Try new actions to discover better rewards
    \item Exploitation: Use known good actions
    \item Balancing these is a key challenge
    \item Common strategies: $\varepsilon$-greedy, UCB, Thompson sampling
\end{itemize}

\textbf{Categories of RL Methods:}
\begin{itemize}[leftmargin=*,noitemsep,topsep=0pt]
    \item Model-based: Learn environment model, then plan
    \item Model-free: Learn directly from experience
    \item Value-based: Learn value function, derive policy
    \item Policy-based: Learn policy directly
\end{itemize}
\end{tcolorbox}

%%%%%%%%%%%%%%%%%%%%%%%%%%%%%%%%%%%%%%%%%%
\begin{tcolorbox}[breakable, title=11.1 Q-Learning and SARSA, colback=rl]
\textbf{When to use:} Use these model-free RL algorithms when the environment model (transition/reward functions) is unknown and must be learned through interaction.

\textbf{Q-Function:} $Q(s,a) = r(s,a) + \gamma \sum_{s'} p(s'|s,a) \max_{a'} Q(s',a')$

\textbf{SARSA Update Rule:}
$Q(s, a) \leftarrow Q(s, a) + \alpha[r + \gamma Q(s', a') - Q(s, a)]$

\textbf{Application Steps:}
\begin{enumerate}[leftmargin=*,noitemsep,topsep=0pt]
    \item From state $s$, choose action $a$ ($\epsilon$-greedy)
    \item Observe reward $r$ and next state $s'$
    \item Choose next action $a'$ ($\epsilon$-greedy)
    \item Update Q-value using formula above
    \item Set $s \leftarrow s'$, $a \leftarrow a'$, continue
\end{enumerate}

\textbf{Q-Learning Update Rule:}
$Q(s, a) \leftarrow Q(s, a) + \alpha[r + \gamma \max_{a'} Q(s', a') - Q(s, a)]$

\textbf{Application Steps:}
\begin{enumerate}[leftmargin=*,noitemsep,topsep=0pt]
    \item From state $s$, choose action $a$ ($\epsilon$-greedy)
    \item Observe reward $r$ and next state $s'$
    \item Update Q-value using formula above (with max operation)
    \item Set $s \leftarrow s'$, continue
\end{enumerate}

\textbf{Key Difference:} SARSA is on-policy (uses actual next action), Q-learning is off-policy (uses best possible next action)

\textbf{Convergence Conditions:}
\begin{itemize}[leftmargin=*,noitemsep,topsep=0pt]
    \item All state-action pairs visited infinitely often
    \item Learning rate schedule: $\sum_t \alpha_t = \infty$ and $\sum_t \alpha_t^2 < \infty$
    \item SARSA requires exploration to decrease over time
\end{itemize}

\textbf{Deep Q-Networks (DQN):}
\begin{itemize}[leftmargin=*,noitemsep,topsep=0pt]
    \item Use neural network to approximate Q-function
    \item Experience replay: Store and reuse past experiences
    \item Target network: Stabilize training with periodic updates
    \item Loss function: $L(w) = \mathbb{E}[(r + \gamma \max_{a'} Q(s',a';w') - Q(s,a;w))^2]$
\end{itemize}
\end{tcolorbox}

%%%%%%%%%%%%%%%%%%%%%%%%%%%%%%%%%%%%%%%%%%
\begin{tcolorbox}[breakable, title=Time Complexity Analysis, colback=white]
\textbf{How to Calculate Time Complexity:}

\textbf{1. Identify basic operations}
\begin{itemize}[leftmargin=*,noitemsep,topsep=0pt]
    \item Arithmetic operations: O(1)
    \item Array/list access: O(1)
    \item Loop iterations: O(n) for n iterations
    \item Nested loops: O(n·m) for outer loop n times, inner loop m times
\end{itemize}

\textbf{2. Count operations in terms of input size}
\begin{itemize}[leftmargin=*,noitemsep,topsep=0pt]
    \item N = number of data points
    \item D = dimensionality of data
    \item K = number of clusters/classes
    \item T = number of iterations
\end{itemize}

\textbf{3. Focus on dominant terms}
\begin{itemize}[leftmargin=*,noitemsep,topsep=0pt]
    \item Ignore constants and lower-order terms
    \item Consider worst-case scenario
\end{itemize}

\textbf{Example Analysis: K-means}
\begin{enumerate}[leftmargin=*,noitemsep,topsep=0pt]
    \item For each iteration (T iterations):
    \begin{itemize}[leftmargin=*,noitemsep]
        \item For each data point (N points):
            \begin{itemize}[leftmargin=*,noitemsep]
                \item Calculate distance to each centroid (K centroids):
                    \begin{itemize}[leftmargin=*,noitemsep]
                        \item Each distance calculation: O(D) operations
                    \end{itemize}
            \end{itemize}
        \item Update K centroids: O(N·D) operations
    \end{itemize}
    \item Total complexity: O(T·N·K·D)
\end{enumerate}

\textbf{Common Algorithm Complexities:}
\begin{itemize}[leftmargin=*,noitemsep,topsep=0pt]
    \item Linear Regression: O(Nd² + d³) for optimal solution
    \item Gradient Descent: O(Ndt) for t iterations
    \item K-means: O(NKdt) for t iterations
    \item Decision Trees: O(dN log N) for training
    \item SVMs: O(N² to N³) depending on implementation
    \item Neural Networks: $O(N_t \cdot \sum_{\ell} n_{\ell} n_{\ell-1})$ for training
    \item PCA: $O(Nd^2 + d^3)$ for computation
    \item HMM Forward-Backward: $O(TK^2)$ for sequence length $T$, $K$ states
\end{itemize}
\end{tcolorbox}

\end{multicols}
\end{document}